%----------------------------------------------------------------------------------------
%	PACKAGES AND OTHER DOCUMENT CONFIGURATIONS
%----------------------------------------------------------------------------------------

\documentclass[10pt,stdletter,dateno,sigleft]{newlfm} % Extra options: 'sigleft' for a left-aligned signature, 'stdletternofrom' to remove the from address, 'letterpaper' for US letter paper - consult the newlfm class manual for more options

%\usepackage{charter} % Use the Charter font for the document text
\usepackage{graphics}


\newsavebox{\Luiuc}\sbox{\Luiuc}{\parbox[b]{1.75in}{\vspace{0.5in}
\includegraphics[width=1.2\linewidth]{/Users/hectorbahamonde/RU/Job_Market/logo.png}}} % Company/institution logo at the top left of the page
\makeletterhead{Uiuc}{\Lheader{\usebox{\Luiuc}}}

\newlfmP{sigsize=50pt} % Slightly decrease the height of the signature field
\newlfmP{addrfromphone} % Print a phone number under the sender's address
\newlfmP{addrfromemail} % Print an email address under the sender's address
\PhrPhone{\texttt{p}} % Customize the "Telephone" text
\PhrEmail{\texttt{e}} % Customize the "E-mail" text

\lthUiuc % Print the company/institution logo

\usepackage{color} % for colors
\usepackage{hyperref} % for links

\hypersetup{
  colorlinks = true,
  urlcolor = blue,
  pdfpagemode = UseNone
}

%----------------------------------------------------------------------------------------
%	YOUR NAME AND CONTACT INFORMATION
%----------------------------------------------------------------------------------------

\namefrom{\vspace{-3cm}Hector Bahamonde} % Name

\addrfrom{
\today\\ % Date
[12pt]
{\bf Hector Bahamonde}\\
Department of Political Science\\
Rutgers University\\
89 George St. \\ % Address
New Brunswick, NJ 08901
}

\phonefrom{(732) 318-9650} % Phone number

\emailfrom{\href{mailto:hector.bahamonde@rutgers.edu}{Hector.Bahamonde@Rutgers.edu} \\ 
\texttt{w}: \href{http://www.hectorbahamonde.com}{www.HectorBahamonde.com}\\
Download last version \href{http://github.com/hbahamonde/Job_Market/raw/master/Bahamonde_Cover_Letter.pdf}{here}} % Link to last version


%----------------------------------------------------------------------------------------
%	ADDRESSEE AND GREETING/CLOSING
%----------------------------------------------------------------------------------------

\greetto{Dear Members of the Search Committee,} % Greeting text
\closeline{\vspace{-1cm}\includegraphics[width=3cm]{/Users/hectorbahamonde/RU/Administracion/signature.pdf} \hspace{-4cm}} % Closing text




%----------------------------------------------------------------------------------------

\begin{document}
\begin{newlfm}

%----------------------------------------------------------------------------------------
%	LETTER CONTENT
%----------------------------------------------------------------------------------------

% introduction
\vspace{-0.3cm}I am writing to apply for the position available in your institution. I am a doctoral candidate in the Department of Political Science at Rutgers University, where I will earn my Ph.D. in June 2017. My dissertation advisors are Robert Kaufman (political science professor, chair), Daniel Kelemen (political science professor and Jean Monnet chair in European union politics), Douglas Blair (political science and economics, professor), and Paul Poast (political science, University of Chicago).

% intro: my research in general (tesis y paper en LAPS)
My {\bf research} explores the economic origins of state capacities as well as the political economy of institutional development and the role of inequality on democratic development. My methods include historical analyses, quantitative methods and experimental designs. Though my research focuses predominantly on Latin America, any of my current and future research projects can be expanded to other developing countries.


% explain job market paper
My {\bf job market} \href{https://github.com/hbahamonde/IncomeTaxAdoption/raw/master/Bahamonde_IncomeTaxAdoption.pdf}{paper} argues that the modernization of the fiscal apparatus was the product of an inter-elite sectoral conflict process between the agricultural and industrial classes. This paper is embedded into a larger {\bf book manuscript} where I analyze how these structural transformations helped states to make institutional investments that lead to the formation of states with higher capacities. I utilize fine-grained historical case study comparisons, sectoral outputs from 1900 to the present, time-series econometric techniques, hazard models, and a novel earthquake dataset that covers sub-national death tolls from 1900 to the present to measure state capacities. This project builds on the fiscal sociology literature and the dual-sector economy model. My overall findings strongly suggest that national industrialization processes challenged the traditional sector by altering the political and military balance, endangering the political monopoly the traditional sector had inherited after the colonial period. This structural transformation and the subsequent institutional investments gave way to the modernization of the state and their respective societies in a way that differs entirely from modernization theory.


% papers
In addition to the book manuscript, I am currently expanding the findings of a series of papers related to vote-buying and vote-selling, using both observational and experimental data in the Americas. 

{\bf Vote-Selling.} With the support of a generous grant I received from the Graduate School and the Department of Political Science, I utilized an experimental design to carry out two experiments in the U.S. out of a series of experiments to be fielded in Latin America for further comparison. I look at price elasticities on vote-selling relative to individual democratic values. Employing a \emph{list experiment} to capture non-biased answers on socially-condemnable/illegal behaviors (like vote-buying), my identification strategy allowed me to observe the ``tipping point'' where individuals prefer the cash over being a democratic citizen. In a separate study that was already fielded in the U.S., I designed a \emph{conjoint experiment} to identify which of Dahl's democratic dimensions should fail to predict individual propensities of vote-selling. Conjoint experiments allow researchers to directly isolate complex multi-dimensional concepts (such as \emph{support for democracy}) and observe which dimension(s) is/are associated to the outcome of interest (vote-selling). Preliminary results show that when the \emph{liberal} component, particularly the \emph{right to associate} fails, individuals are more likely to sell. In this sense, this project is very innovative as it departs from the common strategy of correlating the standard support for democracy question with vote-buying. My identification strategy is more comprehensive as it \emph{decomposes} ``support for democracy'' in several dimensions which are theoretically and substantively relevant (i.e., Dahl's democracy dimensions). 

{\bf Vote-Buying.} In addition to that I have a working \href{https://github.com/hbahamonde/Clientelism_paper/raw/master/Bahamonde_Clientelism_Paper.pdf}{paper} on vote-buying in Brazil ready to submit for publication. The paper starts by recognizing that there is not an agreement on whether parties target groups or individuals. In fact, most of the times, scholars assume that group-targeting and individual-targeting are interchangeable. What seems to be a major problem, however, is that scholars seem to base their decision on their research designs; ethnographers typically study how parties target \emph{individuals} while experimentalist scholars typically look at how parties target districts/municipalities/states (i.e. \emph{groups}). I developed and tested a theory where parties make use of simultaneous segmented targeting techniques. Groups are preferable by brokers when party machines need to secure higher levels of electoral support, relying on the economies of scale and spillover effects that these groups provide. However, individuals are better targets when they are more identifiable, that is when poor individuals are nested in non-poor contexts or when non-poor individuals are nested in poor groups. Interestingly, I find that non-poor individuals are also targeted. The paper uses observational data, a matched design and a short case study (Brazil) to confirm the assumptions made in the models.

% agenda: structure / behavior
As a comparativist and political economist, I believe that advanced methods should be used to keep answering important/big questions. In this sense, my research also has a disciplinary agenda. My scholarly work for example tries to examine classic problems in comparative political development (like state capacities) while at the same time incorporating cutting-edge econometric techniques. Similarly, my working papers and my experiments, which are framed within the behaviorist tradition, are concerned with fundamental questions regarding democratic theory (i.e. ``broken'' democratic values and vote-selling) and political development. Incorporating the U.S. as a \emph{developed} case will allow me to systematically compare it with \emph{developing} Latin American countries. As I feel that it is very important to make this {\bf methods-substance connection}, my efforts are to ensure that I transmit this message across to all of the graduate students I have mentored or had close professional relationships with. 

% teaching: topics
As an instructor, I am interested in {\bf courses} that are carefully designed to answer big questions in comparative politics, not only from a Latin American perspective, but also across the broader spectrum. As a political economist, I mostly focus on economic development, economic history, institutional development and democratic development/underdevelopment. However, I am also interested in political participation, competitive authoritarianisms, welfare politics, party politics and other relevant topics. In order to study these areas, it seems natural to me to incorporate other relevant cases beyond Latin America. Similarly, though I adopt a structuralist epistemology, I have been broadly trained so other approaches could be considered as well. I feel comfortable teaching comparative politics where other approaches such as political anthropology and political psychology could be incorporated too. Finally, I would feel very comfortable teaching {\bf methods courses} such as regression, maximum likelihood estimation, advanced methods (difference in difference, matching, regression discontinuity designs, etc.) and experimental methodology.


More information, syllabi and papers are available on my website: \href{http://www.hectorbahamonde.com}{www.HectorBahamonde.com}. Thank you for considering my application. I look forward to hearing from you.





%----------------------------------------------------------------------------------------

\end{newlfm}
\end{document}
